% ============================================================
\chapter{Martingales --- Discrete Time}
\label{ch:discrete_martingales}
% ============================================================

The term ``martingale'' comes from the world of gambling: it is the strategy of doubling one's bet after each loss, hoping to recover everything at once. Joseph Leo Doob borrowed this word to designate a far deeper mathematical concept: a process whose best prediction of the future, given the past, is the present value. Martingales are the central tool of modern probability theory: they appear in convergence theorems, stochastic integration, mathematical finance (discounted prices are martingales under the risk-neutral measure), and optimal stopping theory.

\section{Definitions}

\begin{definition}[Discrete-Time Martingale]
Let $(\Omega, \mathcal{F}, (\mathcal{F}_n)_{n \geq 0}, \PP)$ be a filtered
probability space. An adapted process $(M_n)_{n \geq 0}$ is a \emph{martingale}
if:
\begin{enumerate}
  \item $\E[\abs{M_n}] < \infty$ for all $n \geq 0$,
  \item $\E[M_{n+1} \mid \mathcal{F}_n] = M_n$ a.s.\ for all $n \geq 0$.
\end{enumerate}
Analogously:
\begin{itemize}
  \item \emph{supermartingale}: $\E[M_{n+1} \mid \mathcal{F}_n] \leq M_n$ a.s.
  \item \emph{submartingale}: $\E[M_{n+1} \mid \mathcal{F}_n] \geq M_n$ a.s.
\end{itemize}
\end{definition}

\begin{remark}
By the tower property, the martingale condition is equivalent to: for all
$m \leq n$, $\E[M_n \mid \mathcal{F}_m] = M_m$ a.s. In particular,
$\E[M_n] = \E[M_0]$ for all $n$.
\end{remark}

\begin{example}
\begin{enumerate}
  \item \textbf{Random walk}: if $(Y_k)_{k \geq 1}$ are i.i.d.\ centred and
    $S_n = \sum_{k=1}^n Y_k$, then $(S_n)$ is a martingale with respect to
    $\mathcal{F}_n = \sigma(Y_1, \ldots, Y_n)$.
  \item \textbf{Martingale transform}: if $(M_n)$ is a martingale and $(H_n)$
    is predictable and bounded, then
    $(H \cdot M)_n = \sum_{k=1}^n H_k (M_k - M_{k-1})$ is a martingale.
  \item \textbf{Exponential martingale}: if $(Y_k)$ are i.i.d.\ and
    $\psi(\theta) = \log \E[e^{\theta Y_1}] < \infty$, then
    $\exp(\theta S_n - n\psi(\theta))$ is a martingale.
  \item \textbf{L\'evy martingale}: if $X \in L^1$ and $\mathcal{F}_n$ is a
    filtration, then $M_n = \E[X \mid \mathcal{F}_n]$ is a martingale.
\end{enumerate}
\end{example}

\section{Martingale Transform}

\begin{definition}[Martingale Transform]
Let $(M_n)_{n \geq 0}$ be a martingale and $(H_n)_{n \geq 1}$ a predictable
process. The \emph{martingale transform} is
\[
  (H \cdot M)_n = \sum_{k=1}^{n} H_k (M_k - M_{k-1}), \quad (H \cdot M)_0 = 0.
\]
\end{definition}

\begin{theorem}[Stability]
If $(M_n)$ is a martingale (resp.\ supermartingale) and $(H_n)$ is predictable
with $0 \leq H_n \leq C$ for a constant $C$, then $(H \cdot M)_n$ is a
martingale (resp.\ supermartingale).
\end{theorem}

\begin{intuition}[title=\textbf{Gambling Strategy}]
The martingale transform models the cumulative gains of a gambler who bets $H_n$
at step $n$ in a fair game $(M_n)$. Predictability means the bet is decided
before observing the outcome. The stability theorem says that no bounded betting
strategy can turn a fair game into a favourable one.
\end{intuition}

\section{Fundamental Inequalities}

\begin{theorem}[Doob's Maximal Inequality]
\label{thm:doob_maximal}
Let $(M_n)_{0 \leq n \leq N}$ be a submartingale. Then for all $\lambda > 0$:
\[
  \lambda \, \PP\Bigl(\max_{0 \leq n \leq N} M_n \geq \lambda\Bigr)
  \leq \E[M_N^+].
\]
\end{theorem}

\begin{theorem}[Doob's $L^p$ Inequality]
\label{thm:doob_Lp}
Let $(M_n)_{0 \leq n \leq N}$ be a martingale and $p > 1$. Then
\[
  \E\Bigl[\max_{0 \leq n \leq N} \abs{M_n}^p\Bigr]
  \leq \Bigl(\frac{p}{p-1}\Bigr)^p \E[\abs{M_N}^p].
\]
\end{theorem}

\begin{theorem}[Doob's $L^2$ Inequality]
\label{thm:doob_L2}
For a square-integrable martingale $(M_n)$:
\[
  \E\Bigl[\max_{0 \leq n \leq N} M_n^2\Bigr] \leq 4 \E[M_N^2].
\]
\end{theorem}

\section{Doob Decomposition}

\begin{theorem}[Doob Decomposition]
\label{thm:doob_decomp}
Every adapted integrable process $(X_n)_{n \geq 0}$ admits a unique decomposition:
\[
  X_n = M_n + A_n,
\]
where $(M_n)$ is a martingale and $(A_n)$ is a predictable process with $A_0 = 0$.
The process $A$ is given by
\[
  A_n = \sum_{k=1}^{n} \bigl(\E[X_k \mid \mathcal{F}_{k-1}] - X_{k-1}\bigr).
\]
If $(X_n)$ is a submartingale, then $(A_n)$ is increasing.
\end{theorem}

\section{Convergence Theorems}

\begin{theorem}[Martingale Convergence --- Doob's Theorem]
\label{thm:doob_convergence}
Let $(M_n)_{n \geq 0}$ be a submartingale with $\sup_n \E[M_n^+] < \infty$.
Then there exists a r.v.\ $M_{\infty} \in L^1$ such that $M_n \to M_{\infty}$
a.s.
\end{theorem}

\begin{proof}[Proof sketch]
We use Doob's \emph{upcrossing lemma}: if $U_N([a,b])$ denotes the number of
upcrossings of $(M_n)_{0 \leq n \leq N}$ from $a$ to $b$ (with $a < b$), then
\[
  (b - a) \E[U_N([a,b])] \leq \E[(M_N - a)^+].
\]
Under $\sup_n \E[M_n^+] < \infty$, we get $\E[U_{\infty}([a,b])] < \infty$
for all rational $a < b$, which implies a.s.\ convergence.
\end{proof}

\begin{warning}[title=\textbf{A.s.\ vs.\ $L^1$ Convergence}]
Doob's theorem gives a.s.\ convergence, but \emph{not} $L^1$ convergence in
general. For $M_n \to M_{\infty}$ in $L^1$, an additional condition (uniform
integrability) is needed.
\end{warning}

\begin{definition}[Uniform Integrability]
A family $(X_i)_{i \in I}$ of r.v.\ is \emph{uniformly integrable} (UI) if
\[
  \lim_{a \to \infty} \sup_{i \in I} \E[\abs{X_i} \mathbf{1}_{\{\abs{X_i} > a\}}] = 0.
\]
\end{definition}

\begin{theorem}[$L^1$ Convergence]
\label{thm:L1_convergence}
Let $(M_n)$ be a martingale. The following are equivalent:
\begin{enumerate}
  \item $(M_n)$ is uniformly integrable,
  \item $M_n$ converges a.s.\ and in $L^1$ to $M_{\infty}$,
  \item there exists $X \in L^1$ such that $M_n = \E[X \mid \mathcal{F}_n]$ for all $n$.
\end{enumerate}
In this case, $M_{\infty} = X$ a.s.\ and the martingale is called \emph{closed}.
\end{theorem}

\begin{corollary}[$L^p$ Convergence for $p > 1$]
If $(M_n)$ is a martingale bounded in $L^p$ for $p > 1$ (i.e.\
$\sup_n \E[\abs{M_n}^p] < \infty$), then $(M_n)$ converges a.s.\ and in $L^p$.
\end{corollary}

\section{Optional Stopping Theorem}

\begin{theorem}[Optional Stopping --- Bounded Version]
\label{thm:ost_bounded}
Let $(M_n)_{n \geq 0}$ be a martingale and $\tau$ a \emph{bounded} stopping time
($\tau \leq N$ a.s.\ for some $N < \infty$). Then
\[
  \E[M_{\tau}] = \E[M_0].
\]
\end{theorem}

\begin{proof}
Write $M_{\tau} = M_0 + \sum_{k=1}^N H_k (M_k - M_{k-1})$ where
$H_k = \mathbf{1}_{\{\tau \geq k\}}$ is predictable (since
$\{\tau \geq k\} = \{\tau \leq k-1\}^c \in \mathcal{F}_{k-1}$). Since the
martingale transform is a martingale, $\E[M_{\tau}] = \E[M_0]$.
\end{proof}

\begin{theorem}[Optional Stopping --- General Version]
\label{thm:ost_general}
Let $(M_n)_{n \geq 0}$ be a uniformly integrable martingale and $\sigma, \tau$
stopping times with $\sigma \leq \tau$ a.s. Then
\[
  \E[M_{\tau} \mid \mathcal{F}_{\sigma}] = M_{\sigma} \quad \text{a.s.}
\]
In particular, $\E[M_{\tau}] = \E[M_0]$.
\end{theorem}

\begin{warning}[title=\textbf{Conditions for Optional Stopping}]
The optional stopping theorem does \emph{not} apply to an arbitrary stopping
time without additional conditions. Classic counterexample: the symmetric random
walk $(S_n)$ with $\tau = \inf\{n : S_n = 1\}$. We have $\E[S_0] = 0$ but
$S_{\tau} = 1$, since $\tau$ is unbounded and $(S_n)$ is not UI.
\end{warning}

\section{Applications}

\subsection{Gambler's Ruin Revisited}

\begin{example}
Random walk $(S_n)$ with $S_0 = i$ and absorbing barriers at $0$ and $N$.
Set $\tau = \inf\{n : S_n \in \{0, N\}\}$.

If $p \neq 1/2$, the exponential martingale $M_n = (q/p)^{S_n}$ is bounded
between $(q/p)^0 = 1$ and $(q/p)^N$, so $\tau$ is a bounded stopping time for
$(M_{n \wedge \tau})$. By optional stopping:
\[
  (q/p)^i = \E[(q/p)^{S_\tau}] = (q/p)^0 \PP_i(\text{ruin})
  + (q/p)^N (1 - \PP_i(\text{ruin})),
\]
whence $\PP_i(\text{ruin}) = \frac{(q/p)^i - (q/p)^N}{1 - (q/p)^N}$.
\end{example}

\subsection{Law of the Iterated Logarithm --- First Step}

\begin{proposition}
If $(M_n)$ is a square-integrable martingale with
$\E[(M_n - M_{n-1})^2 \mid \mathcal{F}_{n-1}] = \sigma^2$ a.s., then
$\langle M \rangle_n = n\sigma^2$ and $M_n^2 - n\sigma^2$ is a martingale.
\end{proposition}

\section{Predictable Quadratic Variation}

\begin{definition}[Predictable Quadratic Variation]
Let $(M_n)$ be a square-integrable martingale. The \emph{predictable quadratic
variation} $(\langle M \rangle_n)$ is the unique predictable increasing process
such that $(M_n^2 - \langle M \rangle_n)$ is a martingale. It is given by:
\[
  \langle M \rangle_n = \sum_{k=1}^{n} \E[(M_k - M_{k-1})^2 \mid \mathcal{F}_{k-1}].
\]
\end{definition}

\begin{theorem}[Convergence via Quadratic Variation]
\label{thm:qv_convergence}
Let $(M_n)$ be a square-integrable martingale. If
$\langle M \rangle_{\infty} = \lim_{n} \langle M \rangle_n < \infty$ a.s.,
then $(M_n)$ converges a.s.\ and in $L^2$.
\end{theorem}

\section{Exercises}

\begin{exercise}
Show that if $(M_n)$ is a martingale and $\varphi$ is convex, then
$(\varphi(M_n))$ is a submartingale (assuming integrability).
\end{exercise}

\begin{exercise}
Let $(X_n)_{n \geq 0}$ be a simple symmetric random walk on $\Z$ starting from
$0$. Show that $X_n^2 - n$ is a martingale.
\end{exercise}

\begin{exercise}
Let $\tau = \inf\{n \geq 0 : \abs{S_n} = a\}$ where $(S_n)$ is the symmetric
random walk. Use the optional stopping theorem to compute $\E[\tau]$.
\end{exercise}

\begin{exercise}
Prove Doob's maximal inequality (Theorem~\ref{thm:doob_maximal}).
\emph{Hint: use $M_N^+ \geq M_{\sigma}^+$ where
$\sigma = \inf\{n : M_n \geq \lambda\} \wedge N$.}
\end{exercise}

\begin{exercise}
Let $(M_n)$ be a positive martingale with $M_0 = 1$. Show that
$\PP(\sup_n M_n \geq a) \leq 1/a$ for all $a > 0$.
\end{exercise}

\begin{exercise}
Let $(M_n)$ be a martingale bounded in $L^2$. Show that it is uniformly
integrable.
\end{exercise}
